\documentclass[../../main.tex]{subfiles}% 注意这里的文件路径不能用 ./main.tex ,否则用latexmk编译子文件会报错
\graphicspath{{\subfix{./image/}}} % 指定图片目录,后续可以直接使用图片文件名
% 注意这里的文件路径不能用 ../../image/ ,否则用latexmk编译子文件会报错

% 例如:
% \begin{figure}[H]
% \centering
% \includegraphics[scale=0.3]{图.png}
% \caption{}
% \label{figure:图}
% \end{figure}
% 注意:上述\label{}一定要放在\caption{}之后,否则引用图片序号会只会显示??.

\begin{document}

\section{Baire定理及其应用}


\begin{theorem}[Baire定理]\label{theorem:Baire定理}
设$E\subset\mathbb{R}^n$是$F_{\sigma}$集，即$E = \bigcup_{k = 1}^{\infty}F_k$，$\,F_k$ ($k = 1,2,\cdots$) 是闭集. 若每个$F_k$皆无内点，则$E$也无内点.
\end{theorem}
\begin{proof}
若$E$有内点，设为$x_0$，则存在$\delta_0>0$，使$\overline{B}(x_0,\delta_0)\subset E$. 因为$F_1$是无内点的，所以必存在$x_1\in B(x_0,\delta_0)$，且有$x_1\notin F_1$. 又因为$F_1$是闭集，所以可以取到$\delta_1(0<\delta_1<1)$，使得
\begin{align*}
\overline{B}(x_1,\delta_1)\cap F_1=\varnothing,
\end{align*}
同时有$\overline{B}(x_1,\delta_1)\subset B(x_0,\delta_0)$. 再从$\overline{B}(x_1,\delta_1)$出发以类似的推理使用于$F_2$，则可得$\overline{B}(x_2,\delta_2)\cap F_2=\varnothing$，同时有$\overline{B}(x_2,\delta_2)\subset B(x_1,\delta_1)$，这里可以要求$0<\delta_2<1/2$. 继续这一过程，可得点列$\{x_k\}$与正数列$\{\delta_k\}$，使得对每个自然数$k$，有
\begin{align*}
\overline{B}(x_k,\delta_k)\subset B(x_{k - 1},\delta_{k - 1}), \quad \overline{B}(x_k,\delta_k)\cap F_k=\varnothing,
\end{align*}
其中$0<\delta_k<1/k$. 由于当$l>k$时，有$x_l\in B(x_k,\delta_k)$，故
\begin{align*}
|x_l - x_k|<\delta_k<\frac{1}{k}.
\end{align*}
这说明$\{x_k\}$是$\mathbb{R}^n$中的基本列（Cauchy列），从而是收敛列，即存在$x\in\mathbb{R}^n$，使得$\lim_{k\to\infty}|x_k - x| = 0$.

此外，从不等式
\begin{align*}
|x - x_k|\leqslant|x - x_l|+|x_l - x_k|<|x - x_l|+\delta_k, \quad l>k
\end{align*}
立即可知（令$l\to\infty$）$|x - x_k|\leqslant\delta_k$. 这说明$x\in\overline{B}(x_k,\delta_k)$，即对一切$k$，$x\notin F_k$. 这与$x\in E$发生矛盾. 

\end{proof}

\begin{example}\label{example:例13}
有理数集$\mathbb{Q}$不是$G_{\delta}$集.
\end{example}
\begin{proof}
事实上，令$\mathbb{Q}=\{r_k: k = 1,2,\cdots\}$，假定$\mathbb{Q}=\bigcap_{i = 1}^{\infty}G_i$，式中$G_i$ ($i = 1,2,\cdots$) 是开集，则有表示式
\begin{align*}
\mathbb{R}=(\mathbb{R}\setminus\mathbb{Q})\cup\mathbb{Q}=\left(\bigcup_{i = 1}^{\infty}G_i^c\right)\cup\left(\bigcup_{k = 1}^{\infty}\{r_k\}\right),
\end{align*}
这里的每个单点集$\{r_k\}$与$G_i^c$皆为闭集，而且从$\overline{G}_i=\mathbb{R}^1$可知每个$G_i^c$是无内点的. 这说明$\mathbb{R}$是可列个无内点之闭集的并集. 从而由Baire定理可知$\mathbb{R}$也无内点，这一矛盾说明$\mathbb{Q}$不是$G_{\delta}$集. 

\end{proof}

\begin{definition}
设$E\subset\mathbb{R}^n$. 若$\overline{E}=\mathbb{R}^n$，则称$E$为$\mathbb{R}^n$中的\textbf{稠密集}；若$\mathring{\overline{E}}=\varnothing$，则称$E$为$\mathbb{R}^n$中的\textbf{无处稠密集}；可数个无处稠密集的并集称为\textbf{贫集}或\textbf{第一纲集}. 不是第一纲集称为\textbf{第二纲集}.
\end{definition}

\begin{example}\label{example:例14}
设$\{G_k\}$是$\mathbb{R}^n$中的稠密开集列，则$G_0 = \bigcap_{k = 1}^{\infty}G_k$在$\mathbb{R}^n$中稠密.
\end{example}
\begin{proof}
只需指出对$\mathbb{R}^n$中任一闭球$\overline{B}=\overline{B}(x,\delta)$，均有$G_0\cap\overline{B}\neq\varnothing$即可. 采用反证法：假定存在闭球$\overline{B}_0=\overline{B}(x_0,\delta_0)$，使得$G_0\cap\overline{B}_0=\varnothing$，则易知
\begin{align*}
\mathbb{R}^n&=(G_0\cap\overline{B}_0)^c = G_0^c\cup(\overline{B}_0)^c,\\
\overline{B}_0&=\mathbb{R}^n\cap\overline{B}_0 = G_0^c\cap\overline{B}_0=\left(\bigcap_{k = 1}^{\infty}G_k\right)^c\cap\overline{B}_0=\bigcup_{k = 1}^{\infty}(G_k^c\cap\overline{B}_0).
\end{align*}
注意到$G_k^c$是无内点的闭集，故由\hyperref[theorem:Baire定理]{Baire定理}可知，$\overline{B}_0$也无内点，矛盾.

\end{proof}

\begin{example}\label{example:例15}
设$f_k\in C(\mathbb{R}^n)$ ($k = 1,2,\cdots$). 若$\lim_{k\to\infty}f_k(x)=f(x)$ ($x\in\mathbb{R}^n$)，则$f(x)$的不连续点集为第一纲集.
\end{example}
\begin{proof}
注意到$f(x)$的连续点集的表示，只需指出(\refexa{example:例12})
\[
\left(G\left(\frac{1}{m}\right)\right)^c \quad \left(G\left(\frac{1}{m}\right)=\bigcup_{k = 1}^{\infty}\mathring{E}_k\left(\frac{1}{m}\right)\right)
\]
是第一纲集. 对$\varepsilon>0$，令
\[
F_k(\varepsilon)=\bigcap_{i = 1}^{\infty}\left\{x\in\mathbb{R}^n: |f_k(x)-f_{k + i}(x)|\leqslant\varepsilon\right\},
\]
易知$\mathbb{R}^n=\bigcup_{k = 1}^{\infty}F_k(\varepsilon)$，$F_k(\varepsilon)\subset E_k(\varepsilon)$，从而有
\[
\mathring{F}_k(\varepsilon)\subset\mathring{E}_k(\varepsilon)\subset G(\varepsilon), \quad \bigcup_{k = 1}^{\infty}\mathring{F}_k(\varepsilon)\subset G(\varepsilon).
\]
由此知
\begin{align*}
[G(\varepsilon)]^c&=\mathbb{R}^n\setminus G(\varepsilon)\subset\mathbb{R}^n\setminus\bigcup_{k = 1}^{\infty}\mathring{F}_k(\varepsilon)\\
&=\bigcup_{k = 1}^{\infty}F_k(\varepsilon)\setminus\bigcup_{k = 1}^{\infty}\mathring{F}_k(\varepsilon)\subset\bigcup_{k = 1}^{\infty}[F_k(\varepsilon)\setminus\mathring{F}_k(\varepsilon)]=\bigcup_{k = 1}^{\infty}\partial F_k(\varepsilon).
\end{align*}
因为$F_k(\varepsilon)$是闭集，所以$\partial F_k(\varepsilon)$是无处稠密集. 这说明$(G(\varepsilon))^c$是第一纲集. 

\end{proof}

\begin{example}\label{example:例16}
设$f\in C([0,1])$，且令
\[
f_1'(x)=f(x), f_2'(x)=f_1(x), \cdots, f_n'(x)=f_{n - 1}(x), \cdots.
\]
若对每一个$x\in[0,1]$，都存在自然数$k$，使得$f_k(x)=0$，则$f(x)\equiv0$.
\end{example}
\begin{proof}
只需指出$f(x)$在$[0,1]$中的一个稠密集上为$0$即可. 对此，我们在$[0,1]$中任取一个闭子区间$I$，并记
\[
F_k=\{x\in I: f_k(x)=0\} \quad (k = 1,2,\cdots).
\]
显然，每个$F_k$都是闭集，且$I=\bigcup_{k = 1}^{\infty}F_k$. 根据\hyperref[theorem:Baire定理]{Baire定理}可知，存在$F_{k_0}$，它包含一个区间$(\alpha,\beta)$. 因为在$(\alpha,\beta)$上$f_{k_0}(x)=0$，所以$f(x)=0$，$x\in(\alpha,\beta)$. 注意到$(\alpha,\beta)\subset I$，即得所证. 

\end{proof}

\begin{definition}[Baire第一函数类]
称区间$I$上连续函数列的极限函数$f(x)$的全体为\textbf{Baire第一函数类}，记为$f\in B_1(I)$.
\end{definition}

\begin{theorem}\label{theorem:Baire第一函数类的性质1}
若$f_n\in B_1(\mathbb{R})$，且$f_n(x)$在$\mathbb{R}$上一致收敛于$f(x)$，则$f\in B_1(\mathbb{R})$.
\end{theorem}
\begin{proof}
事实上，由题设知，对任意$k\in\mathbb{N}$，存在$n_k\in\mathbb{N}$，使得
\begin{align*}
|f_{n_k}(x)-f(x)|<1/2^{k + 1} \quad (x\in\mathbb{R}).
\end{align*}
这里不妨认定$n_1 < n_2 < \cdots < n_k < \cdots$. 考查$\sum_{k = 1}^{\infty}[f_{n_{k + 1}}(x)-f_{n_k}(x)]$，因为我们有
\begin{align*}
|f_{n_{k + 1}}(x)-f_{n_k}(x)| &\leqslant |f_{n_{k + 1}}(x)-f(x)|+|f_{n_k}(x)-f(x)| \\
&<\frac{1}{2^{k + 2}}+\frac{1}{2^{k + 1}}<\frac{1}{2^{k}} \quad (x\in\mathbb{R}),
\end{align*}
所以$g(x)\stackrel{\text{def}}{=}\sum_{k = 1}^{\infty}[f_{n_{k + 1}}(x)-f_{n_k}(x)]\in B_1(\mathbb{R})$. 显然有$g(x)=f(x)-f_{n_1}(x)$，即$f(x)=g(x)+f_{n_1}(x)$. 证毕.

\end{proof}


\end{document}